\input{preamble}

\title{Section 6.2 --- Lines --- Answer Key}

%argument 1: m
%argument 2: x1
%argument 3: y1
%argument 4: b
%argument 5: standard form
\newcommand{\threeqn}[5]{%
a. $y #3 = #1\left(x #2\right)$\\
b. $y=#1 x #4$\\
c. $#5$}

\begin{document}
\maketitle

\section*{Equations of Lines}
For the lines described in each exercise:
\begin{clist}
\item Write an equation of the line in point-slope form.
\item Rewrite the equation in slope-intercept form.
\item Rewrite the equation in standard form.
\end{clist}
\begin{smenum}
\mitemxxx
{\threeqn{-\frac{3}{5}}{-2}{-5}{+\frac{31}{5}}{3 x +5 y = 31}}
{\threeqn{3}{+2}{-\frac{3}{8}}{+\frac{51}{8}}{24 x -8 y = -51}}
%\mitemxx
{\threeqn{\frac{3}{5}}{-2}{+4}{-\frac{26}{5}}{3 x -5 y = 26}}
\mitemxxx
{\threeqn{-2}{-3}{+2}{+4}{2 x + y = 4}}
%\mitemxx
{\threeqn{-\frac{2}{3}}{+1}{+1}{-\frac{5}{3}}{2 x +3 y = -5}}
{\threeqn{3}{-3}{-3}{-6}{3 x - y = 6}}
\end{smenum}
For each line described below, write an equation in whatever form is easiest.
\begin{smenum}
\mitemxxxx
{$2x+5y=-13$}
{$3x-2y=-7$}
%\mitemxx
{$9x+8y=20$}
{$2x-3y=-19$}
\mitemxxxx
{$x=-2$}
{$y=8$}
%\mitemxx
{$y=2$}
{$x=5$}
\end{smenum}

\section*{Systems of Equations}
Solve each system below by graphing both lines on the same axes and finding the point of intersection. You can use the graph paper on the other side of this sheet, or a separate sheet of graph paper.
\begin{smenum}
\mitemxx
{\begin{GridSquareFilled}{5}
\lineslopeintercept(-5,-5)(5,5){3}{1/2}
\lineslopeintercept(-5,-5)(5,5){7}{-3/2}
\filldraw (2,4) circle (\dotsize);
\end{GridSquareFilled}
$x=2$, $y=4$}
{\begin{GridSquareFilled}{5}
\lineslopeintercept(-5,-5)(5,5){4}{-2/3}
\lineslopeintercept(-5,-5)(5,5){4}{-2/3}
\end{GridSquareFilled}
Infinitely many solutions}
\end{smenum}
Solve each system below by substitution or elimination, whichever is easier.
\begin{smenum}
\maitemxxxx
{x=-20,\ y=0}
{x=2,\ y=-3}
%\maitemxx
{x=8,\ y=-3}
{x=0,\ y=6}
\end{smenum}

\section{Distance From a Line}
For each point and line given below, give the distance from the point to the line.
\begin{smenum}
\mitemxxxx
{$2\sqrt{5}$}
{$\tfrac{7}{5}\sqrt{10}$}
{$1$}
{$5$}
\end{smenum}
\end{document}